\subsubsection{Interfaces con el Usuario}
% Contenido pendiente de redacción.
La interacción con el comandante o coordinador se realizará mediante los dispositivos descritos en la sección~2.1.2 y deberá cumplir los siguientes requerimientos:

\begin{itemize}
  \item \textbf{Instalación.} La aplicación deberá entregarse como un ejecutable para Windows~11 que funcione sin instalar previamente Python ni otros programas~\cite{pmp2026}. La instalación deberá incluir \textcolor{red}{el dataset de escenarios}, y su procedimiento se describirá en el manual de instalación y uso.
  \item \textbf{Ventana principal.} La aplicación deberá presentar una ventana principal, diseñada para una resolución de 1920~$\times$~1080 píxeles, con el mapa georreferenciado del sector de Aguas Claras en la mayor parte de la pantalla y un panel lateral con los datos de la consulta, sus resultados y los mensajes de estado. El mapa deberá poder acercarse con la rueda del ratón y desplazarse al arrastrarlo.
  \item \textbf{Idioma.} La interfaz deberá estar disponible en español e inglés, y el idioma deberá poder cambiarse desde el panel de configuraciones avanzadas descrito en la sección~2.1.6.
  \item \textbf{Ubicación y perímetro del incendio.} El usuario deberá poder señalar la ubicación del incendio con un clic sobre el mapa o trazar con clics sucesivos el perímetro observado, e indicar el tramo que corresponde al borde activo y el territorio ya quemado. La aplicación deberá permitir deshacer el último punto, borrar el trazo y modificarlo para actualizar la consulta cuando haya nuevas observaciones, y deberá advertir cuando la ubicación quede fuera del área cubierta \textcolor{red}{por el dataset}.
  \item \textbf{Condiciones meteorológicas.} Si la consulta a Open-Meteo falla o no hay conexión a internet, la aplicación deberá mostrar un formulario para ingresar con el teclado la velocidad del viento en km/h, la dirección de procedencia del viento en grados desde el norte (de $0^{\circ}$ a $360^{\circ}$), la presencia de lluvia y su cantidad en milímetros, sin descartar la ubicación ya señalada. El formulario deberá rechazar los valores fuera de rango e indicar el rango válido.
  \item \textbf{Resultados.} Al completar la consulta, la aplicación deberá mostrar sobre el mapa la estimación de la propagación y, en el panel lateral, las condiciones meteorológicas utilizadas con su origen: Open-Meteo o ingreso manual. Desde ese panel, el usuario deberá poder abrir ala animación.
  \item \textbf{Facilidad de uso.} En la evaluación con al menos tres bomberos, la interfaz deberá obtener un puntaje promedio mínimo de 68 en el cuestionario de usabilidad SUS, y al menos el 80~\% de las tareas deberá completarse sin ayuda~\cite{pmp2026}.
\end{itemize}